%% ─────────────────────────────────────────────────────────
%%  Chapter 3 — Requirement Engineering
%%  Formatting: longtables have caption, endfirsthead, endhead,
%%  endfoot, endlastfoot. Column widths are within 14.5 cm.
%% ─────────────────────────────────────────────────────────

\section{Feasibility Study}

\subsection{Technical Feasibility}
The system is built entirely on accessible, open-source technologies. The embedding model
(\texttt{all-MiniLM-L6-v2}) runs locally without GPU requirements, making it deployable on
standard academic workstations. The vector store (ChromaDB in-memory mode) requires no
database server. The generative LLM is accessed via the Groq API, which provides free-tier
access to Llama 3.3~70B. The Streamlit frontend requires only a Python environment.
All dependencies are installable via \texttt{pip}. The entire system runs on a single machine
with 8\,GB RAM, confirming full technical feasibility.

\subsection{Operational Feasibility}
The system is designed for two user classes: a \textit{patient} role accessing public
guidelines and a research/demo mode for evaluating attack scenarios. The interface is a
web-based Streamlit application requiring no installation from the end-user's perspective.
Attack and defence states are toggled through sidebar controls, making the system operable
without technical expertise. For researchers, the evaluation script
(\texttt{scripts/evaluate.py}) automates the full 100-query benchmark with a single command.

\subsection{Economic Feasibility}
All components are open-source or free-tier:
\begin{itemize}
    \item Groq API: Free tier with 14{,}400 tokens/minute limit, sufficient for evaluation.
    \item Sentence-Transformers: Open-source (Apache 2.0 licence).
    \item ChromaDB: Open-source (Apache 2.0 licence).
    \item Streamlit: Open-source (Apache 2.0 licence).
\end{itemize}
The total monetary cost of the project is \$0, confirming economic feasibility.

\subsection{Legal and Ethical Feasibility}
All documents in the knowledge base are sourced from \textbf{publicly available clinical
guidelines} published by the NHS, CDC, and WHO. No real patient data is used. The attack
scenarios are simulated on synthetic adversarial documents created specifically for research.
The project explicitly does not process, store, or transmit any personally identifiable
information (PII) or protected health information (PHI), and therefore does not fall under
HIPAA's Privacy Rule. All research activities comply with standard academic ethics guidelines
for AI security research.

%% ─────────────────────────────────────────────────────────
\section{Requirement Elicitation and Analysis}

\subsection{Stakeholders}
\begin{itemize}
    \item \textbf{Primary User (Patient):} Accesses only public-classified documents.
          Receives grounded, safety-scored answers to healthcare queries.
    \item \textbf{Security Researcher:} Uses the Attack Demo mode to observe the
          three-condition comparison (clean, poisoned, defended) and analyse per-layer
          defence effectiveness.
    \item \textbf{System Administrator:} Uses the sidebar controls to inject poisoned
          documents, reset the database, and monitor clean vs.\ poisoned document counts.
\end{itemize}

\subsection{Methods Used for Requirement Elicitation}

Requirements were elicited through three complementary methods, chosen to match the
research-oriented nature of the project:
\begin{enumerate}
    \item \textbf{Literature Analysis:} The functional security requirements were derived
          directly from the attack taxonomy documented in PoisonedRAG (Zou et al., 2025),
          Greshake et al.\ (2023), and the OWASP LLM Top~10 (2025). Each identified attack
          vector (knowledge poisoning, indirect injection, cross-tenant leakage) corresponds
          to at least one implemented defence layer.
    \item \textbf{Threat Modelling using MITRE ATLAS:} Non-functional security requirements
          (NFR-02, NFR-03, NFR-07) were derived by mapping the project's threat surface to
          the MITRE ATLAS adversarial taxonomy (AML.T0043, AML.T0051, AML.T0020).
    \item \textbf{Iterative Prototyping:} UI and pipeline requirements evolved through three
          implementation cycles. The attack demo three-column layout was added after the
          first prototype revealed that side-by-side comparison was essential for
          demonstrating the Implicit Trust Paradox. No direct user interviews were conducted
          given the research nature of the system.
\end{enumerate}

\subsection{Use Cases}
\begin{enumerate}
    \item \textbf{UC-01:} Patient submits a medical query and receives a risk-scored,
          grounded answer.
    \item \textbf{UC-02:} Researcher runs the attack demo to observe how a poisoned RAG
          produces incorrect medical advice.
    \item \textbf{UC-03:} Researcher activates the defence and observes CRIR reduction in
          the defended column.
    \item \textbf{UC-04:} Administrator injects poisoned documents via the sidebar and
          verifies the DB count updates.
    \item \textbf{UC-05:} Administrator resets the database to the clean baseline state.
    \item \textbf{UC-06:} Researcher expands the Defence Logs to inspect which documents
          were retrieved, blocked, and why.
\end{enumerate}
\begin{figure}[H]
\centering
\makebox[\textwidth]{ % <-- THIS FIXES TRUE CENTERING

\begin{tikzpicture}[
  actor/.style={font=\small, text centered, minimum width=1.5cm},
  usecase/.style={ellipse, draw=black!70, fill=blue!6, text centered,
                  minimum width=3.2cm, minimum height=0.75cm, font=\footnotesize},
  assoc/.style={thin, black!60}
]
%% System boundary
\draw[draw=black!50, fill=gray!3, rounded corners=4pt]
  (1.2,-0.3) rectangle (9.8,-8.5);
\node[font=\small\bfseries] at (5.5, -0.7) {MediCare Secure RAG System};

%% Use case ellipses (inside boundary)
\node[usecase] at (5.5,-1.8)  (uc1) {UC-01: Submit medical query};
\node[usecase] at (5.5,-3.0)  (uc2) {UC-02: Run attack demo};
\node[usecase] at (5.5,-4.2)  (uc3) {UC-03: Activate defence};
\node[usecase] at (5.5,-5.4)  (uc4) {UC-04: Inject poisoned docs};
\node[usecase] at (5.5,-6.6)  (uc5) {UC-05: Reset database};
\node[usecase] at (5.5,-7.8)  (uc6) {UC-06: Inspect Defence Logs};

%% Actors (outside boundary)
\node[actor] at (-1.0,-1.8) (patient) {Patient};
\draw[fill=black!70] (-1.0,-1.0) circle (0.2cm);
\draw[black!70] (-1.0,-1.2) -- (-1.0,-2.0);
\draw[black!70] (-1.4,-1.5) -- (-0.6,-1.5);
\draw[black!70] (-1.0,-2.0) -- (-1.3,-2.5);
\draw[black!70] (-1.0,-2.0) -- (-0.7,-2.5);

\node[actor] at (-1.0,-5.0) (researcher) {Researcher};
\draw[fill=black!70] (-1.0,-4.2) circle (0.2cm);
\draw[black!70] (-1.0,-4.4) -- (-1.0,-5.2);
\draw[black!70] (-1.4,-4.7) -- (-0.6,-4.7);
\draw[black!70] (-1.0,-5.2) -- (-1.3,-5.7);
\draw[black!70] (-1.0,-5.2) -- (-0.7,-5.7);

\node[actor] at (11.5,-6.0) (admin) {Administrator};
\draw[fill=black!70] (11.5,-5.2) circle (0.2cm);
\draw[black!70] (11.5,-5.4) -- (11.5,-6.2);
\draw[black!70] (11.1,-5.7) -- (11.9,-5.7);
\draw[black!70] (11.5,-6.2) -- (11.2,-6.7);
\draw[black!70] (11.5,-6.2) -- (11.8,-6.7);

%% Associations
\draw[assoc] (patient)    -- (uc1);
\draw[assoc] (researcher) -- (uc2);
\draw[assoc] (researcher) -- (uc3);
\draw[assoc] (researcher) -- (uc6);
\draw[assoc] (admin)      -- (uc4);
\draw[assoc] (admin)      -- (uc5);
\draw[assoc] (admin)      -- (uc6);

\end{tikzpicture}

} % <-- end makebox

\caption{UML use case diagram showing system boundary, three actor types (Patient, Researcher, Administrator), and six identified use cases with actor associations.}
\label{fig:usecase}
\end{figure}
\section{Software Requirement Specification (SRS)}

\subsection{Functional Requirements}

\setlength{\LTleft}{\fill}
\setlength{\LTright}{\fill}

\renewcommand{\arraystretch}{1.05} % slightly compact

\begin{longtable}{|p{1.2cm}|p{3.8cm}|p{7.0cm}|}
\caption{Functional Requirements of MediCare Secure RAG.}
\label{tab:fr} \\
\hline
\rowcolor{gray!15}
\textbf{ID} & \textbf{Requirement} & \textbf{Description} \\
\hline
\endfirsthead

\multicolumn{3}{l}{\small\textit{Table \thetable{} continued from previous page}} \\[2pt]
\hline
\rowcolor{gray!15}
\textbf{ID} & \textbf{Requirement} & \textbf{Description} \\
\hline
\endhead

\hline
\multicolumn{3}{r}{\small\textit{Continue in next page}} \\  % <-- your required wording
\endfoot

\hline
\endlastfoot

FR-01 & Document Ingestion       & System shall load clean JSONL documents into the
                                    ChromaDB in-memory collection on startup. \\
\hline
FR-02 & Adversarial Injection    & System shall support injection of 27 poisoned documents
                                    via a sidebar button without resetting the clean corpus. \\
\hline
FR-03 & Semantic Retrieval       & System shall retrieve top-5 semantically similar
                                    documents using cosine similarity over sentence embeddings. \\
\hline
FR-04 & Defence Layer Execution  & System shall pass retrieved documents through Layers~2,
                                    1, and~5 sequentially (in that order), blocking
                                    adversarial documents at each stage. \\
\hline
FR-05 & Hardened Generation      & System shall pass only approved documents to the LLM
                                    using a positional-delimiter secure prompt template
                                    (Layer~3). \\
\hline
FR-06 & Output Monitoring        & System shall scan generated responses for 8 defined
                                    danger patterns using compiled regex (Layer~4). \\
\hline
FR-07 & CRIR Computation         & System shall compute the binary CRIR score (0 or 1)
                                    and retrieval precision metric for every query. \\
\hline
FR-08 & Risk Scoring             & System shall compute a composite risk score (0--100)
                                    from CRIR, blocked count, groundedness, and validation
                                    issues. \\
\hline
FR-09 & Attack Demo              & System shall execute a sequential three-condition
                                    comparison (clean, poisoned, defended) for any
                                    user-supplied query. \\
\hline
FR-10 & Database Reset           & System shall reset the ChromaDB collection to the clean
                                    baseline on administrator request. \\
\hline

\end{longtable}

\subsection{Non-Functional Requirements}

\setlength{\LTleft}{\fill}
\setlength{\LTright}{\fill}

\renewcommand{\arraystretch}{1.05}

\begin{longtable}{|p{1.2cm}|p{3.2cm}|p{7.5cm}|}
\caption{Non-Functional Requirements of MediCare Secure RAG.}
\label{tab:nfr} \\
\hline
\rowcolor{gray!15}
\textbf{ID} & \textbf{Category} & \textbf{Requirement} \\
\hline
\endfirsthead

\multicolumn{3}{l}{\small\textit{Table \thetable{} continued from previous page}} \\[2pt]
\hline
\rowcolor{gray!15}
\textbf{ID} & \textbf{Category} & \textbf{Requirement} \\
\hline
\endhead

\hline
\multicolumn{3}{r}{\small\textit{Continue in next page}} \\
\endfoot

\hline
\endlastfoot

NFR-01 & Performance       & Response latency shall not exceed 8 seconds per query under
                             normal network conditions. \\
\hline
NFR-02 & Accuracy          & Clean DB CRIR score shall be 0.0 across all evaluation queries
                             (zero false positives from the defence pipeline). \\
\hline
NFR-03 & Security          & The system shall not expose API keys, system prompts, or
                             internal database contents in any response. \\
\hline
NFR-04 & Usability         & The Streamlit interface shall be operable without any
                             installation steps beyond \texttt{pip install} and
                             \texttt{streamlit run}. \\
\hline
NFR-05 & Reproducibility   & The evaluation script shall produce identical summary statistics
                             on the same hardware given the same JSONL inputs. \\
\hline
NFR-06 & Maintainability   & All defence layers shall be implemented in isolated,
                             independently testable Python functions. \\
\hline
NFR-07 & Compliance        & The knowledge base shall contain exclusively publicly available
                             clinical guidelines with no PII or PHI. \\
\hline

\end{longtable}
%% ─────────────────────────────────────────────────────────
\section{Technology Stack}

\setlength{\LTleft}{\fill}
\setlength{\LTright}{\fill}

\renewcommand{\arraystretch}{1.05}

\begin{longtable}{|p{3.0cm}|p{4.2cm}|p{4.8cm}|}
\caption{Technology Stack and Primary Security Role of Each Component.}
\label{tab:techstack} \\
\hline
\rowcolor{gray!15}
\textbf{Component} & \textbf{Technology} & \textbf{Primary Security Role} \\
\hline
\endfirsthead

\multicolumn{3}{l}{\small\textit{Table \thetable{} continued from previous page}} \\[2pt]
\hline
\rowcolor{gray!15}
\textbf{Component} & \textbf{Technology} & \textbf{Primary Security Role} \\
\hline
\endhead

\hline
\multicolumn{3}{r}{\small\textit{Continue in next page}} \\
\endfoot

\hline
\endlastfoot

Vector Store      & ChromaDB (in-memory)                          & Secure storage and metadata-filtered retrieval \\
\hline
Embeddings        & Sentence-Transformers (\texttt{all-MiniLM-L6-v2}) & Semantic mapping and anomaly detection \\
\hline
Generative LLM    & Groq-hosted Llama 3.3 70B                    & Final reasoning and medical synthesis \\
\hline
Frontend          & Streamlit                                     & UI-based attack simulation and monitoring \\
\hline
Language          & Python 3.10+                                  & Core orchestration and all defence layers \\
\hline
Environment       & \texttt{python-dotenv}                        & Secure API key management via \texttt{.env} \\
\hline
HTTP Client       & Groq Python SDK                               & LLM API communication \\
\hline

\end{longtable}